%!TEX program = tectonic
\documentclass[12pt,a4paper]{report}

\usepackage[spanish]{babel}
\usepackage[T1]{fontenc}
\usepackage[utf8]{inputenc}
\usepackage[a4paper,margin=2.5cm]{geometry}
\usepackage{graphicx}
\usepackage{float}
\usepackage{booktabs}
\usepackage{longtable}
\usepackage{array}
\usepackage{enumitem}
\usepackage{xcolor}
\usepackage{tikz}
\usepackage{microtype}
\usepackage[hidelinks]{hyperref}

\usetikzlibrary{arrows.meta,calc,positioning,shapes.geometric,shapes.multipart}

\setlength{\parskip}{0.5em}
\setlength{\parindent}{0pt}
\setlength{\emergencystretch}{3em}
\raggedbottom
\renewcommand{\arraystretch}{1.15}
\renewcommand{\contentsname}{Tabla de contenido}

\newcommand{\dercasversion}{1.0}
\newcommand{\dercasdate}{15 de octubre de 2026}
\newcommand{\dercasestado}{Borrador para verificación: incluye defectos conocidos y trabajo pendiente.}
\newcommand{\dercassistema}{FoodLoop: red social de recetas}
\newcommand{\cucampo}[2]{\textbf{#1} & #2 \\[0.35em]}

\begin{document}

\begin{titlepage}
  \centering
  \vspace*{2.5cm}
  {\Huge\bfseries FOODLOOP\par}
  \vspace{0.4cm}
  {\Large Red social de recetas\par}
  \vspace{1.8cm}
  {\large Documento de Especificaciones, Requerimientos y\\
  Criterios de Aceptación de Software - DERCAS\par}
  \vspace{1.8cm}
  {\large Departamento de Desarrollo de Sistemas, DTI\par}
  \vspace{0.5cm}
  {\large octubre 2026\par}
  \vfill
  {\normalsize Versión \dercasversion\par}
  \vspace{0.4cm}
  {\small \dercasestado\par}
\end{titlepage}

\chapter*{Historial de revisiones}
\addcontentsline{toc}{chapter}{Historial de revisiones}

\begin{center}
\begin{tabular}{p{3cm}p{2cm}p{8.5cm}}
  \toprule
  Fecha & Versión & Descripción \\
  \midrule
  15 de octubre de 2026 & 1.0 & Primera versión del DERCAS a partir del código real, la API implementada, los diagramas Markdown y los defectos conocidos. Reestructurado según el modelo de referencia (12 secciones). \\
  \bottomrule
\end{tabular}
\end{center}

\tableofcontents

\chapter{Definiciones, Acrónimos y Abreviaturas}
\label{ch:definiciones}

\begin{center}
\begin{longtable}{p{4.2cm}p{9.3cm}}
  \toprule
  Término & Concepto \\
  \midrule
  \endfirsthead
  \toprule
  Término & Concepto \\
  \midrule
  \endhead
  Receta & Documento creado por un usuario con título, descripción, ingredientes, instrucciones, categoría, etiquetas y visibilidad. Puede tener una imagen en GridFS. \\
  Visibilidad & Atributo \texttt{publica} de una receta. Si es verdadero, cualquier usuario autenticado puede verla; si es falso, solo su propietario. \\
  Propietario & Usuario que creó una receta. Es el único que puede editarla, eliminarla, cambiar su visibilidad y gestionar su imagen. \\
  Imagen & Archivo binario JPEG, PNG, WebP o GIF guardado en el bucket GridFS \texttt{uploads}. La receta conserva el identificador del archivo. \\
  Interacción social & Acción alternable sobre una receta: me gusta, guardar o compartir. El reporte escribe un motivo y marca la receta. \\
  Comentario & Texto asociado a una receta, con posibilidad de respuesta de primer nivel. Su borrado es lógico. \\
  Reseña & Calificación de 1 a 5 con texto opcional. Un usuario puede dejar una sola reseña por receta. \\
  Seguimiento & Relación dirigida entre dos usuarios: uno sigue al otro. Es alternable. \\
  Feed & Lista paginada de actividad propia y de los usuarios seguidos. \\
  Campana & Vista provisional del feed usada como notificaciones. No persiste el estado de lectura. \\
  Actividad & Registro de una acción social sobre una receta o usuario, usado para construir el feed. \\
  Reporte & Marca de contenido inadecuado con un motivo. No hay moderación automática ni cuenta administradora. \\
  DERCAS & Documento de Especificación de Requerimientos para el Análisis y Diseño del Sistema. \\
  API & Interfaz de programación de aplicaciones. \\
  REST & Estilo de arquitectura de servicios basado en recursos y verbos HTTP. \\
  CRUD & Crear, leer, actualizar y eliminar. \\
  JWT & Token web JSON usado para autenticar peticiones. \\
  GridFS & Mecanismo de MongoDB para almacenar archivos binarios en fragmentos. \\
  MIME & Tipo de medio declarado de un archivo (\texttt{image/jpeg}, etc.). \\
  VPS & Servidor privado virtual usado para producción. \\
  TLS & Seguridad de la capa de transporte para HTTPS. \\
  CORS & Intercambio de recursos de origen cruzado. \\
  RNF & Requerimiento no funcional. \\
  \texttt{ObjectId} & Identificador de documentos MongoDB. \\
  \texttt{CastError} & Error de Mongoose al convertir un identificador inválido. \\
  \bottomrule
\end{longtable}
\end{center}

\chapter{Descripción del sistema solicitado}
\label{ch:descripcion}

El proyecto parte de un recetario personal y se transforma en una red social de recetas. Se solicitó el desarrollo e implementación de un sistema para el registro y la gestión de recetas que permita:

\begin{itemize}[leftmargin=*]
  \item[a.] El registro de usuarios y la autenticación mediante un token JWT de 30 días.
  \item[b.] La creación, consulta, edición, eliminación y visibilidad pública o privada de recetas.
  \item[c.] El almacenamiento de una imagen por receta en el bucket GridFS \texttt{uploads}.
  \item[d.] Las interacciones sociales sobre cada receta: me gusta, guardar, compartir y reportar.
  \item[e.] Los comentarios, las respuestas de primer nivel y las reseñas con promedio.
  \item[f.] El seguimiento entre usuarios, el feed de actividad y una campana de notificaciones.
\end{itemize}

El sistema usa un cliente Angular, una API Express, MongoDB Atlas para documentos y GridFS para imágenes. Implementa la red social básica y conserva los defectos conocidos que se documentan en los capítulos \ref{ch:procesos}, \ref{ch:casos-uso} y \ref{ch:planificacion}.

\chapter{Definición de actores}
\label{ch:actores}

\begin{center}
\begin{longtable}{p{4cm}p{9.5cm}}
  \toprule
  Nombre del actor & Descripción \\
  \midrule
  \endfirsthead
  \toprule
  Nombre del actor & Descripción \\
  \midrule
  \endhead
  Visitante & Persona sin sesión. Puede registrarse, iniciar sesión y, si conoce la URL, ver una imagen pública. No puede llamar a rutas protegidas. \\
  Usuario registrado & Persona autenticada. Puede crear y gestionar sus recetas, interactuar socialmente, comentar, reseñar, seguir usuarios y consultar el feed. \\
  Propietario de receta & Usuario registrado con derechos adicionales sobre una receta concreta: editarla, eliminarla, gestionar su imagen y cambiar su visibilidad. La propiedad se determina por receta, no por un rol global. \\
  Sistema de almacenamiento & MongoDB Atlas y GridFS como actores externos de persistencia para documentos e imágenes. \\
  \bottomrule
\end{longtable}
\end{center}

No existe una cuenta administradora. La distinción entre usuario común y propietario depende de la propiedad de cada receta.

\chapter{Procesos a sistematizar}
\label{ch:procesos}

Los procesos utilizan los identificadores P01 a P06 y se describen con los campos del modelo de referencia: nombre, descripción, actores, entradas, salidas, especificaciones de funcionamiento y restricciones adicionales. Las referencias detalladas por endpoints están en \path{docs/referencia/api.md}, los diagramas de flujo en el capítulo \ref{ch:flujo} y los casos de uso en el capítulo \ref{ch:casos-uso}.

\section{P01: Autenticación y sesión}

\begin{center}
\begin{longtable}{p{3.6cm}p{9.9cm}}
  \toprule
  Campo & Contenido \\
  \midrule
  \endfirsthead
  \toprule
  Campo & Contenido \\
  \midrule
  \endhead
  Nombre & P01: Autenticación y sesión. \\
  Descripción & Registro, inicio de sesión y cierre local de sesión. \\
  Actores & Visitante y usuario registrado. \\
  Entradas & Nombre, correo y contraseña. \\
  Salidas & Identificador, nombre, correo y token JWT. \\
  Especificaciones de funcionamiento & \path{POST /api/auth/register} crea el usuario con hash bcrypt; \path{POST /api/auth/login} compara el hash y emite un token de 30 días. El cierre de sesión elimina el token del navegador. Alcance: acceso a la aplicación y a las rutas protegidas. El correo debe ser único y la contraseña debe tener al menos 8 caracteres. \\
  Restricciones adicionales & Límite de 10 intentos cada 15 minutos y saneado de nombre y correo en el registro. Respuestas \texttt{201}, \texttt{200}, \texttt{400}, \texttt{429} y \texttt{500}. \\
  \bottomrule
\end{longtable}
\end{center}

\section{P02: Gestión de recetas y visibilidad}

\begin{center}
\begin{longtable}{p{3.6cm}p{9.9cm}}
  \toprule
  Campo & Contenido \\
  \midrule
  \endfirsthead
  \toprule
  Campo & Contenido \\
  \midrule
  \endhead
  Nombre & P02: Gestión de recetas y visibilidad. \\
  Descripción & Creación, consulta, edición, eliminación y cambio de visibilidad de recetas. \\
  Actores & Usuario registrado y propietario. \\
  Entradas & Título, descripción, ingredientes, instrucciones, categoría y visibilidad. \\
  Salidas & Receta creada, lista filtrada, detalle permitido o confirmación de borrado. \\
  Especificaciones de funcionamiento & \path{POST /api/recipes}, \path{GET /api/recipes}, \path{GET /api/recipes?scope=public}, \path{GET /api/recipes/:id}, \path{PUT /api/recipes/:id} y \path{DELETE /api/recipes/:id}. El servidor asigna el propietario y acepta únicamente 11 campos editables. Las recetas nacen privadas; el detalle permite recetas propias o públicas; la edición y el borrado exigen propiedad. Alcance: tablero personal, recetas públicas y detalle. \\
  Restricciones adicionales & Lista de campos editables, validadores Mongoose, límite de creación y saneado parcial. No hay búsqueda ni paginación en el servidor. Respuestas \texttt{201}, \texttt{200}, \texttt{400}, \texttt{404}, \texttt{429} y \texttt{500}. \\
  \bottomrule
\end{longtable}
\end{center}

\section{P03: Gestión de imágenes}

\begin{center}
\begin{longtable}{p{3.6cm}p{9.9cm}}
  \toprule
  Campo & Contenido \\
  \midrule
  \endfirsthead
  \toprule
  Campo & Contenido \\
  \midrule
  \endhead
  Nombre & P03: Gestión de imágenes. \\
  Descripción & Selección, optimización, subida, reemplazo, eliminación y visualización de imágenes. \\
  Actores & Propietario para la gestión; cualquier persona con la URL para la visualización. \\
  Entradas & Archivo de imagen e identificador de receta. \\
  Salidas & Identificador GridFS asociado o receta sin imagen. \\
  Especificaciones de funcionamiento & \path{POST /api/recipes/:id/image} con \texttt{multipart/form-data}; \path{PUT /api/recipes/:id} con imagen vacía; \path{GET /api/recipes/image/:fileId} público. Solamente se conserva la imagen vigente; la anterior se borra al reemplazar, quitar o eliminar la receta. Alcance: formulario de edición, tarjetas, detalle y enlace directo. \\
  Restricciones adicionales & MIME de imagen, 25 MB en el navegador, 5 MB en el servidor y propiedad de la receta. No se verifican los bytes reales. Respuestas \texttt{200}, \texttt{400}, \texttt{404} y \texttt{500}. \\
  \bottomrule
\end{longtable}
\end{center}

\section{P04: Interacciones sociales}

\begin{center}
\begin{longtable}{p{3.6cm}p{9.9cm}}
  \toprule
  Campo & Contenido \\
  \midrule
  \endfirsthead
  \toprule
  Campo & Contenido \\
  \midrule
  \endhead
  Nombre & P04: Interacciones sociales. \\
  Descripción & Alternar me gusta, guardar, compartir y reportar recetas. \\
  Actores & Usuario registrado. \\
  Entradas & Identificador de receta y, para reportes, motivo. \\
  Salidas & Estado alternado y contadores actualizados. \\
  Especificaciones de funcionamiento & \path{POST /api/recipes/:id/like}, \path{/save}, \path{/share} y \path{/report}. El me gusta crea o elimina una fila en \texttt{likes}; guardar usa \texttt{saves}. Las acciones son alternables; dar me gusta a la receta propia responde \texttt{400}; compartir registra actividad. Alcance: tarjetas y modal del tablero. \\
  Restricciones adicionales & Límite social de 30 acciones por minuto, límite de creación para reportes y saneado del motivo. Respuestas \texttt{200}, \texttt{201} (reporte), \texttt{400} (me gusta propio o reporte duplicado), \texttt{404}, \texttt{429} y \texttt{500}. \\
  \bottomrule
\end{longtable}
\end{center}

\section{P05: Comentarios y reseñas}

\begin{center}
\begin{longtable}{p{3.6cm}p{9.9cm}}
  \toprule
  Campo & Contenido \\
  \midrule
  \endfirsthead
  \toprule
  Campo & Contenido \\
  \midrule
  \endhead
  Nombre & P05: Comentarios y reseñas. \\
  Descripción & Conversación y calificación de una receta. \\
  Actores & Usuario registrado y autor del contenido. \\
  Entradas & Texto, identificador de receta, comentario padre y puntuación. \\
  Salidas & Comentario o reseña creado, promedio actualizado y contadores. \\
  Especificaciones de funcionamiento & \path{POST /api/comments}, \path{GET /api/comments/receta/:recetaId}, \path{PUT /api/comments/:id}, \path{DELETE /api/comments/:id}, y las rutas equivalentes en \path{/api/reviews}. Un usuario puede publicar una sola reseña por receta; el promedio se redondea a un decimal; el borrado de comentarios es lógico. Alcance: modal de detalle. \\
  Restricciones adicionales & Texto obligatorio, límites de longitud, puntuación de 1 a 5 y control de autor. La lectura solamente expone un nivel de respuestas y la edición o borrado aún no tiene botones. Respuestas \texttt{201}, \texttt{200}, \texttt{400}, \texttt{404}, \texttt{429} y \texttt{500}. \\
  \bottomrule
\end{longtable}
\end{center}

\section{P06: Seguimiento, feed y notificaciones}

\begin{center}
\begin{longtable}{p{3.6cm}p{9.9cm}}
  \toprule
  Campo & Contenido \\
  \midrule
  \endfirsthead
  \toprule
  Campo & Contenido \\
  \midrule
  \endhead
  Nombre & P06: Seguimiento, feed y notificaciones. \\
  Descripción & Relaciones de seguimiento, descubrimiento mediante feed, perfil y campana provisional. \\
  Actores & Usuario registrado. \\
  Entradas & Identificador del usuario seguido, página y límite del feed. \\
  Salidas & Estado de seguimiento, listas, actividades y contador provisional. \\
  Especificaciones de funcionamiento & \path{POST /api/follows}, listas y comprobación de seguimiento, \path{GET /api/activity/feed} y \path{GET /api/activity/user/:userId}. No existe colección de notificaciones. El seguimiento es alternable; el feed incluye actividades dirigidas a seguidos y al propio usuario. Alcance: rutas \texttt{/feed} y \texttt{/profile/:username}, además de la campana del tablero. \\
  Restricciones adicionales & Prohibición de seguirse a sí mismo, pareja única de seguimiento, paginación del feed y trabajo pendiente para resolver perfiles por nombre y persistir lectura. Respuestas \texttt{200}, \texttt{400}, \texttt{404}, \texttt{429} y \texttt{500}. \\
  \bottomrule
\end{longtable}
\end{center}

\chapter{Diagrama de flujo de los procesos}
\label{ch:flujo}

Los diagramas se generaron con Mermaid y las fuentes editables están en \path{docs/diagramas/flujo}. Se presentan ordenados desde el flujo general hasta cada proceso.

\section{Flujo general del sistema}

\begin{figure}[H]
  \centering
  \includegraphics[width=\textwidth,height=0.85\textheight,keepaspectratio]{build/flujo/flujo-sistema-general.pdf}
  \caption{Flujo general del sistema.}
  \label{fig:flujo-general}
\end{figure}

\section{P01: Autenticación y sesión}

\begin{figure}[H]
  \centering
  \includegraphics[width=\textwidth,height=0.85\textheight,keepaspectratio]{build/flujo/flujo-autenticacion-y-sesion.pdf}
  \caption{Flujo de autenticación y sesión.}
  \label{fig:flujo-autenticacion}
\end{figure}

\section{P02: Recetas y visibilidad}

\begin{figure}[H]
  \centering
  \includegraphics[width=\textwidth,height=0.85\textheight,keepaspectratio]{build/flujo/flujo-recetas-y-visibilidad.pdf}
  \caption{Flujo de gestión de recetas y visibilidad.}
  \label{fig:flujo-recetas}
\end{figure}

\section{P03: Imágenes}

\begin{figure}[H]
  \centering
  \includegraphics[width=\textwidth,height=0.85\textheight,keepaspectratio]{build/flujo/flujo-imagenes.pdf}
  \caption{Flujo del ciclo de vida de la imagen.}
  \label{fig:flujo-imagenes}
\end{figure}

\section{P04: Interacciones sociales}

\begin{figure}[H]
  \centering
  \includegraphics[width=\textwidth,height=0.85\textheight,keepaspectratio]{build/flujo/flujo-interacciones-sociales.pdf}
  \caption{Flujo de las interacciones sociales.}
  \label{fig:flujo-interacciones}
\end{figure}

\section{P05: Comentarios y reseñas}

\begin{figure}[H]
  \centering
  \includegraphics[width=\textwidth,height=0.85\textheight,keepaspectratio]{build/flujo/flujo-comentarios-y-resenas.pdf}
  \caption{Flujo de comentarios y reseñas.}
  \label{fig:flujo-comentarios}
\end{figure}

\section{P06: Seguimiento, feed y notificaciones}

\begin{figure}[H]
  \centering
  \includegraphics[width=\textwidth,height=0.85\textheight,keepaspectratio]{build/flujo/flujo-seguimiento-feed-y-notificaciones.pdf}
  \caption{Flujo de seguimiento, feed y campana provisional.}
  \label{fig:flujo-seguimiento}
\end{figure}

\chapter{Casos de uso}
\label{ch:casos-uso}

\section{Autenticación y sesión}

\begin{figure}[H]
  \centering
  \includegraphics[width=\textwidth,height=0.55\textheight,keepaspectratio]{build/casos-uso/casos-uso-autenticacion-y-sesion.pdf}
  \caption{Casos de uso de autenticación y sesión.}
  \label{fig:cu-autenticacion}
\end{figure}

\subsection{CU-01: Registrarse}
\begin{center}
\begin{tabular}{p{3.4cm}p{9.9cm}}
  \toprule
  \cucampo{Descripción}{Crear una cuenta y recibir un token.}
  \cucampo{Actores primarios}{Visitante.}
  \cucampo{Contexto}{P01 Autenticación y sesión.}
  \cucampo{Condiciones previas}{El correo no debe estar registrado.}
  \cucampo{Flujo básico}{El visitante envía el formulario. El sistema sanea los campos, comprueba que el correo no exista, aplica el hash bcrypt, crea el usuario y devuelve el token.}
  \cucampo{Requerimientos especiales}{Contraseña de al menos 8 caracteres y correo único.}
  \cucampo{Salida}{Usuario creado con \texttt{201} e ingreso al tablero.}
  \cucampo{Excepciones}{Correo ya registrado (\texttt{400}); contraseña corta (\texttt{500} por validación); exceso de intentos (\texttt{429}).}
  \bottomrule
\end{tabular}
\end{center}

\subsection{CU-02: Iniciar sesión}
\begin{center}
\begin{tabular}{p{3.4cm}p{9.9cm}}
  \toprule
  \cucampo{Descripción}{Autenticarse y recibir un token.}
  \cucampo{Actores primarios}{Visitante.}
  \cucampo{Contexto}{P01 Autenticación y sesión.}
  \cucampo{Condiciones previas}{Tener una cuenta.}
  \cucampo{Flujo básico}{El visitante envía correo y contraseña. El sistema busca al usuario, compara el hash con \texttt{bcrypt.compare} y emite un JWT de 30 días.}
  \cucampo{Requerimientos especiales}{Límite de 10 intentos cada 15 minutos.}
  \cucampo{Salida}{Sesión iniciada con \texttt{200} y token guardado localmente.}
  \cucampo{Excepciones}{Usuario inexistente (\texttt{400}); contraseña incorrecta (\texttt{400}); exceso de intentos (\texttt{429}).}
  \bottomrule
\end{tabular}
\end{center}

\subsection{CU-03: Cerrar sesión}
\begin{center}
\begin{tabular}{p{3.4cm}p{9.9cm}}
  \toprule
  \cucampo{Descripción}{Terminar la sesión en el navegador.}
  \cucampo{Actores primarios}{Usuario registrado.}
  \cucampo{Contexto}{P01 Autenticación y sesión.}
  \cucampo{Condiciones previas}{Haber iniciado sesión.}
  \cucampo{Flujo básico}{El usuario pulsa cerrar sesión; el cliente elimina el token local y vuelve a \path{/login}.}
  \cucampo{Requerimientos especiales}{No invalida el token en el servidor.}
  \cucampo{Salida}{Sesión cerrada localmente.}
  \cucampo{Excepciones}{Un token copiado sigue siendo válido hasta expirar.}
  \bottomrule
\end{tabular}
\end{center}

\section{Recetas y visibilidad}

\begin{figure}[H]
  \centering
  \includegraphics[width=\textwidth,height=0.55\textheight,keepaspectratio]{build/casos-uso/casos-uso-recetas-y-visibilidad.pdf}
  \caption{Casos de uso de recetas y visibilidad.}
  \label{fig:cu-recetas}
\end{figure}

\subsection{CU-04: Crear receta}
\begin{center}
\begin{tabular}{p{3.4cm}p{9.9cm}}
  \toprule
  \cucampo{Descripción}{Guardar una receta propia.}
  \cucampo{Actores primarios}{Usuario registrado.}
  \cucampo{Contexto}{P02 Gestión de recetas y visibilidad.}
  \cucampo{Condiciones previas}{Sesión válida.}
  \cucampo{Flujo básico}{El usuario envía el formulario. El sistema conserva solo los campos editables, asigna el propietario y devuelve \texttt{201}; la receta nace privada.}
  \cucampo{Requerimientos especiales}{Límite de 10 creaciones por minuto.}
  \cucampo{Salida}{Receta creada y visible en el tablero.}
  \cucampo{Excepciones}{Exceso de creaciones (\texttt{429}); error de validación (\texttt{500}).}
  \bottomrule
\end{tabular}
\end{center}

\subsection{CU-05: Listar y ver recetas}
\begin{center}
\begin{tabular}{p{3.4cm}p{9.9cm}}
  \toprule
  \cucampo{Descripción}{Consultar recetas propias o públicas.}
  \cucampo{Actores primarios}{Usuario registrado.}
  \cucampo{Contexto}{P02 Gestión de recetas y visibilidad.}
  \cucampo{Condiciones previas}{Sesión válida.}
  \cucampo{Flujo básico}{El tablero solicita las recetas propias o, con \texttt{scope=public}, las públicas; el detalle permite la receta propia o una pública.}
  \cucampo{Requerimientos especiales}{Sin paginación ni búsqueda en el servidor.}
  \cucampo{Salida}{Lista y detalle mostrados.}
  \cucampo{Excepciones}{Receta privada ajena (\texttt{404}); identificador inválido puede producir \texttt{500}.}
  \bottomrule
\end{tabular}
\end{center}

\subsection{CU-06: Editar, eliminar y cambiar visibilidad}
\begin{center}
\begin{tabular}{p{3.4cm}p{9.9cm}}
  \toprule
  \cucampo{Descripción}{Actualizar, borrar o publicar una receta.}
  \cucampo{Actores primarios}{Propietario.}
  \cucampo{Contexto}{P02 Gestión de recetas y visibilidad.}
  \cucampo{Condiciones previas}{Ser propietario de la receta.}
  \cucampo{Flujo básico}{El propietario guarda cambios, confirma el borrado o alterna la visibilidad. El sistema valida la propiedad, aplica la lista de campos editables y ejecuta los validadores.}
  \cucampo{Requerimientos especiales}{El borrado físico deja huérfanos comentarios, reseñas y actividad.}
  \cucampo{Salida}{Receta actualizada, eliminada o publicada.}
  \cucampo{Excepciones}{Sin campos editables (\texttt{400}); receta ajena o inexistente (\texttt{404}).}
  \bottomrule
\end{tabular}
\end{center}

\section{Imágenes}

\begin{figure}[H]
  \centering
  \includegraphics[width=\textwidth,height=0.55\textheight,keepaspectratio]{build/casos-uso/casos-uso-imagenes.pdf}
  \caption{Casos de uso de imágenes.}
  \label{fig:cu-imagenes}
\end{figure}

\subsection{CU-07: Subir imagen}
\begin{center}
\begin{tabular}{p{3.4cm}p{9.9cm}}
  \toprule
  \cucampo{Descripción}{Asociar una imagen optimizada a una receta existente.}
  \cucampo{Actores primarios}{Propietario.}
  \cucampo{Contexto}{P03 Gestión de imágenes.}
  \cucampo{Condiciones previas}{Receta propia creada.}
  \cucampo{Flujo básico}{El usuario elige un archivo; el navegador valida y optimiza a 1600 px JPEG 0.8; se envía como \texttt{multipart/form-data}; el servidor guarda el nuevo identificador y borra la imagen anterior.}
  \cucampo{Requerimientos especiales}{MIME de imagen, 25 MB en el navegador y 5 MB en el servidor.}
  \cucampo{Salida}{Receta con imagen actualizada (\texttt{200}).}
  \cucampo{Excepciones}{Tipo no permitido o tamaño excedido (\texttt{400}); receta ajena (\texttt{404}).}
  \bottomrule
\end{tabular}
\end{center}

\subsection{CU-08: Reemplazar o quitar imagen}
\begin{center}
\begin{tabular}{p{3.4cm}p{9.9cm}}
  \toprule
  \cucampo{Descripción}{Mantener una sola imagen vigente o dejar la receta sin imagen.}
  \cucampo{Actores primarios}{Propietario.}
  \cucampo{Contexto}{P03 Gestión de imágenes.}
  \cucampo{Condiciones previas}{Receta propia.}
  \cucampo{Flujo básico}{El propietario elige otra imagen o elimina la actual; el sistema registra el nuevo identificador o cadena vacía y elimina el archivo anterior de GridFS.}
  \cucampo{Requerimientos especiales}{Solo una imagen vigente por receta.}
  \cucampo{Salida}{Imagen reemplazada o eliminada.}
  \cucampo{Excepciones}{Un fallo de subida conserva el identificador anterior.}
  \bottomrule
\end{tabular}
\end{center}

\subsection{CU-09: Ver imagen}
\begin{center}
\begin{tabular}{p{3.4cm}p{9.9cm}}
  \toprule
  \cucampo{Descripción}{Mostrar la imagen en tarjetas, detalle o enlace.}
  \cucampo{Actores primarios}{Cualquier persona con la URL pública.}
  \cucampo{Contexto}{P03 Gestión de imágenes.}
  \cucampo{Condiciones previas}{Conocer el identificador del archivo.}
  \cucampo{Flujo básico}{El sistema valida el identificador, lee GridFS, devuelve el tipo de contenido y aplica caché inmutable.}
  \cucampo{Requerimientos especiales}{Lectura pública sin cabecera de autorización.}
  \cucampo{Salida}{Imagen servida con \texttt{200}.}
  \cucampo{Excepciones}{Identificador inválido (\texttt{400}); archivo inexistente (\texttt{404}).}
  \bottomrule
\end{tabular}
\end{center}

\section{Interacciones sociales}

\begin{figure}[H]
  \centering
  \includegraphics[width=\textwidth,height=0.55\textheight,keepaspectratio]{build/casos-uso/casos-uso-interacciones-sociales.pdf}
  \caption{Casos de uso de interacciones sociales.}
  \label{fig:cu-interacciones}
\end{figure}

\subsection{CU-10: Dar o quitar me gusta}
\begin{center}
\begin{tabular}{p{3.4cm}p{9.9cm}}
  \toprule
  \cucampo{Descripción}{Expresar aprobación y mantener el contador.}
  \cucampo{Actores primarios}{Usuario registrado.}
  \cucampo{Contexto}{P04 Interacciones sociales.}
  \cucampo{Condiciones previas}{Sesión válida.}
  \cucampo{Flujo básico}{El usuario pulsa el corazón; el sistema alterna el estado, sincroniza la receta y los favoritos del usuario y registra actividad para el autor.}
  \cucampo{Requerimientos especiales}{Límite social de 30 acciones por minuto.}
  \cucampo{Salida}{Estado alternado con \texttt{liked} y \texttt{likesCount}.}
  \cucampo{Excepciones}{Receta inexistente (\texttt{404}); la interfaz evita el me gusta propio, pero la API no lo impide.}
  \bottomrule
\end{tabular}
\end{center}

\subsection{CU-11: Guardar o quitar guardado}
\begin{center}
\begin{tabular}{p{3.4cm}p{9.9cm}}
  \toprule
  \cucampo{Descripción}{Conservar una receta para consulta posterior.}
  \cucampo{Actores primarios}{Usuario registrado.}
  \cucampo{Contexto}{P04 Interacciones sociales.}
  \cucampo{Condiciones previas}{Sesión válida.}
  \cucampo{Flujo básico}{El usuario pulsa guardar; el sistema alterna el arreglo \texttt{guardados} y su contador.}
  \cucampo{Requerimientos especiales}{No sincroniza el arreglo \texttt{favoritos} del usuario.}
  \cucampo{Salida}{Estado alternado con \texttt{saved} y \texttt{guardadosCount}.}
  \cucampo{Excepciones}{Receta inexistente (\texttt{404}).}
  \bottomrule
\end{tabular}
\end{center}

\subsection{CU-12: Compartir}
\begin{center}
\begin{tabular}{p{3.4cm}p{9.9cm}}
  \toprule
  \cucampo{Descripción}{Aumentar el contador de compartidos y registrar actividad.}
  \cucampo{Actores primarios}{Usuario registrado.}
  \cucampo{Contexto}{P04 Interacciones sociales.}
  \cucampo{Condiciones previas}{Sesión válida.}
  \cucampo{Flujo básico esperado}{El usuario pulsa compartir; el sistema incrementa el contador y crea la actividad \texttt{compartio\_receta}.}
  \cucampo{Requerimientos especiales}{Defecto conocido: el tipo no está en la enumeración de actividad.}
  \cucampo{Salida}{Contador incrementado.}
  \cucampo{Excepciones}{Puede devolver \texttt{500} después de incrementar el contador.}
  \bottomrule
\end{tabular}
\end{center}

\subsection{CU-13: Reportar}
\begin{center}
\begin{tabular}{p{3.4cm}p{9.9cm}}
  \toprule
  \cucampo{Descripción}{Informar contenido inapropiado.}
  \cucampo{Actores primarios}{Usuario registrado.}
  \cucampo{Contexto}{P04 Interacciones sociales.}
  \cucampo{Condiciones previas}{Sesión válida.}
  \cucampo{Flujo básico}{El usuario envía un motivo; el sistema lo sanea, guarda el reporte y marca la receta como reportada.}
  \cucampo{Requerimientos especiales}{Admite reportes repetidos; no hay moderación.}
  \cucampo{Salida}{Receta marcada como reportada.}
  \cucampo{Excepciones}{Receta inexistente (\texttt{404}).}
  \bottomrule
\end{tabular}
\end{center}

\section{Comentarios y reseñas}

\begin{figure}[H]
  \centering
  \includegraphics[width=\textwidth,height=0.55\textheight,keepaspectratio]{build/casos-uso/casos-uso-comentarios-y-resenas.pdf}
  \caption{Casos de uso de comentarios y reseñas.}
  \label{fig:cu-comentarios}
\end{figure}

\subsection{CU-14: Leer comentarios y reseñas}
\begin{center}
\begin{tabular}{p{3.4cm}p{9.9cm}}
  \toprule
  \cucampo{Descripción}{Ver la conversación y la valoración de una receta.}
  \cucampo{Actores primarios}{Usuario registrado.}
  \cucampo{Contexto}{P05 Comentarios y reseñas.}
  \cucampo{Condiciones previas}{Sesión válida.}
  \cucampo{Flujo básico}{Al abrir el detalle, el sistema pide los comentarios principales con un nivel de respuestas y las reseñas ordenadas.}
  \cucampo{Requerimientos especiales}{Solo un nivel de respuestas.}
  \cucampo{Salida}{Comentarios, reseñas y promedio mostrados.}
  \cucampo{Excepciones}{Errores de lectura muestran mensajes genéricos.}
  \bottomrule
\end{tabular}
\end{center}

\subsection{CU-15: Comentar y responder}
\begin{center}
\begin{tabular}{p{3.4cm}p{9.9cm}}
  \toprule
  \cucampo{Descripción}{Participar en la conversación de una receta.}
  \cucampo{Actores primarios}{Usuario registrado.}
  \cucampo{Contexto}{P05 Comentarios y reseñas.}
  \cucampo{Condiciones previas}{Sesión válida y texto no vacío.}
  \cucampo{Flujo básico}{El usuario publica un comentario o respuesta; el sistema lo guarda, incrementa el contador, crea actividades y lo muestra en el modal.}
  \cucampo{Requerimientos especiales}{Límite de 10 creaciones por minuto; borrado lógico.}
  \cucampo{Salida}{Comentario creado con \texttt{201}.}
  \cucampo{Excepciones}{Texto vacío o excesivo; no se comprueba que la receta exista o sea visible.}
  \bottomrule
\end{tabular}
\end{center}

\subsection{CU-16: Reseñar}
\begin{center}
\begin{tabular}{p{3.4cm}p{9.9cm}}
  \toprule
  \cucampo{Descripción}{Calificar una receta de 1 a 5 estrellas.}
  \cucampo{Actores primarios}{Usuario registrado.}
  \cucampo{Contexto}{P05 Comentarios y reseñas.}
  \cucampo{Condiciones previas}{Sesión válida y puntuación distinta de cero.}
  \cucampo{Flujo básico}{El usuario envía la reseña; el sistema la guarda, recalcula el promedio y el total, crea actividad y la muestra.}
  \cucampo{Requerimientos especiales}{Una sola reseña por usuario y receta; promedio redondeado a un decimal.}
  \cucampo{Salida}{Reseña creada con \texttt{201} y promedio actualizado.}
  \cucampo{Excepciones}{Segunda reseña del mismo usuario (\texttt{400}); puntuación inválida puede producir \texttt{500}.}
  \bottomrule
\end{tabular}
\end{center}

\subsection{CU-17: Editar o borrar contenido propio}
\begin{center}
\begin{tabular}{p{3.4cm}p{9.9cm}}
  \toprule
  \cucampo{Descripción}{Corregir o retirar el comentario o reseña propio.}
  \cucampo{Actores primarios}{Autor del contenido.}
  \cucampo{Contexto}{P05 Comentarios y reseñas.}
  \cucampo{Condiciones previas}{Ser autor.}
  \cucampo{Flujo básico}{El autor actualiza el texto o aplica el borrado lógico; el sistema recalcula métricas.}
  \cucampo{Requerimientos especiales}{Existe en la API, pero actualmente no tiene botones en la interfaz.}
  \cucampo{Salida}{Comentario o reseña actualizado o retirado.}
  \cucampo{Excepciones}{Control de autor; recurso inexistente (\texttt{404}).}
  \bottomrule
\end{tabular}
\end{center}

\section{Seguimiento, feed y notificaciones}

\begin{figure}[H]
  \centering
  \includegraphics[width=\textwidth,height=0.55\textheight,keepaspectratio]{build/casos-uso/casos-uso-seguimiento-feed-y-notificaciones.pdf}
  \caption{Casos de uso de seguimiento, feed y notificaciones.}
  \label{fig:cu-seguimiento}
\end{figure}

\subsection{CU-18: Seguir o dejar de seguir}
\begin{center}
\begin{tabular}{p{3.4cm}p{9.9cm}}
  \toprule
  \cucampo{Descripción}{Suscribirse o cancelar la suscripción a otro usuario.}
  \cucampo{Actores primarios}{Usuario registrado.}
  \cucampo{Contexto}{P06 Seguimiento, feed y notificaciones.}
  \cucampo{Condiciones previas}{Sesión válida y no ser el mismo usuario.}
  \cucampo{Flujo básico}{El usuario pulsa seguir; el sistema alterna el documento \texttt{Follow}, actualiza ambos contadores y registra actividad para el usuario seguido.}
  \cucampo{Requerimientos especiales}{Pareja única de seguimiento; límite de 10 creaciones por minuto.}
  \cucampo{Salida}{Estado alternado con \texttt{isFollowing}.}
  \cucampo{Excepciones}{Seguirse a sí mismo (\texttt{400}).}
  \bottomrule
\end{tabular}
\end{center}

\subsection{CU-19: Ver seguidores, seguidos, feed y perfil}
\begin{center}
\begin{tabular}{p{3.4cm}p{9.9cm}}
  \toprule
  \cucampo{Descripción}{Explorar relaciones y actividad.}
  \cucampo{Actores primarios}{Usuario registrado.}
  \cucampo{Contexto}{P06 Seguimiento, feed y notificaciones.}
  \cucampo{Condiciones previas}{Sesión válida.}
  \cucampo{Flujo básico}{El sistema devuelve seguidores, seguidos, actividades paginadas y actividad por usuario.}
  \cucampo{Requerimientos especiales}{Defecto conocido: \path{/profile/:username} envía un nombre donde el backend espera un identificador.}
  \cucampo{Salida}{Listas y feed mostrados.}
  \cucampo{Excepciones}{El perfil no resuelve correctamente usuario, recetas ni seguimiento.}
  \bottomrule
\end{tabular}
\end{center}

\subsection{CU-20: Consultar campana}
\begin{center}
\begin{tabular}{p{3.4cm}p{9.9cm}}
  \toprule
  \cucampo{Descripción}{Revisar la actividad reciente.}
  \cucampo{Actores primarios}{Usuario registrado.}
  \cucampo{Contexto}{P06 Seguimiento, feed y notificaciones.}
  \cucampo{Condiciones previas}{Sesión válida.}
  \cucampo{Flujo básico}{La campana reutiliza el feed, descarta las actividades propias y cuenta todas como no leídas.}
  \cucampo{Requerimientos especiales}{Provisional; marcar como leída solo modifica la memoria local.}
  \cucampo{Salida}{Contador y lista provisional.}
  \cucampo{Excepciones}{Pendiente modelo, endpoints y persistencia (Fase 7).}
  \bottomrule
\end{tabular}
\end{center}

\chapter{Diagrama de Secuencia}
\label{ch:secuencia}

Los diagramas de secuencia muestran la interacción entre el cliente Angular, la API Express, MongoDB y GridFS para los procesos principales. Las fuentes editables están en \path{docs/diagramas/secuencia}.

\section{P01: Registro e inicio de sesión}

\begin{figure}[H]
  \centering
  \includegraphics[width=\textwidth,height=0.85\textheight,keepaspectratio]{build/secuencia/secuencia-registro-e-inicio-sesion.pdf}
  \caption{Secuencia de registro, inicio de sesión y cierre local.}
  \label{fig:seq-auth}
\end{figure}

\section{P02 y P03: Crear receta y gestionar imagen}

\begin{figure}[H]
  \centering
  \includegraphics[width=\textwidth,height=0.85\textheight,keepaspectratio]{build/secuencia/secuencia-crear-receta-y-gestionar-imagen.pdf}
  \caption{Secuencia de creación de receta y gestión de su imagen.}
  \label{fig:seq-receta-imagen}
\end{figure}

\section{P05: Comentar y reseñar}

\begin{figure}[H]
  \centering
  \includegraphics[width=\textwidth,height=0.85\textheight,keepaspectratio]{build/secuencia/secuencia-comentar-y-resenar.pdf}
  \caption{Secuencia de comentarios y reseñas.}
  \label{fig:seq-comentarios}
\end{figure}

\section{P06: Seguir usuarios y consultar el feed}

\begin{figure}[H]
  \centering
  \includegraphics[width=\textwidth,height=0.85\textheight,keepaspectratio]{build/secuencia/secuencia-seguir-usuarios-y-feed.pdf}
  \caption{Secuencia de seguimiento, feed y campana provisional.}
  \label{fig:seq-feed}
\end{figure}

\chapter{Diagramas de estado}
\label{ch:estados}

Los diagramas de estado describen el ciclo de vida de los recursos con estados propios. Las fuentes editables están en \path{docs/diagramas/estado}.

\section{Ciclo de vida de la receta}

\begin{figure}[H]
  \centering
  \includegraphics[width=\textwidth]{build/estado/estado-receta.pdf}
  \caption{Estados de una receta: visibilidad, reporte y eliminación.}
  \label{fig:estado-receta}
\end{figure}

La receta nace privada. La visibilidad y el reporte son atributos del documento, y el borrado es físico: no hay borrado lógico ni cascada.

\begin{itemize}[leftmargin=*]
  \item \texttt{reportado} es un booleano independiente; una receta reportada no cambia de visibilidad ni deja de ser consultable.
  \item No existe estado de moderación ni cuenta administradora.
  \item Al eliminar, se borra la receta y su imagen GridFS, pero quedan huérfanos los comentarios, reseñas y actividades asociados.
\end{itemize}

\section{Ciclo de vida de la imagen}

\begin{figure}[H]
  \centering
  \includegraphics[width=\textwidth]{build/estado/estado-imagen.pdf}
  \caption{Estados de la imagen de una receta.}
  \label{fig:estado-imagen}
\end{figure}

Solo se conserva una imagen vigente por receta. El archivo anterior se elimina de GridFS al reemplazar, quitar o borrar la receta.

\begin{itemize}[leftmargin=*]
  \item La lectura \path{GET /api/recipes/image/:fileId} es pública y usa caché inmutable.
  \item La validación real de bytes (magic bytes) está pendiente; solo se comprueba el tipo MIME declarado y un máximo de 5 MB en el servidor.
\end{itemize}

\section{Estado de la campana de notificaciones}

\begin{figure}[H]
  \centering
  \includegraphics[width=\textwidth]{build/estado/estado-campana-notificaciones.pdf}
  \caption{Estado provisional de la campana de notificaciones.}
  \label{fig:estado-campana}
\end{figure}

La campana es provisional: reutiliza el feed y no persiste el estado de lectura.

\begin{itemize}[leftmargin=*]
  \item El contador de no leídas es la longitud de la lista, no un estado real.
  \item \texttt{markAsRead()} solo modifica el arreglo en memoria del navegador.
  \item El modelo \texttt{Notification} con \texttt{destinatario}, \texttt{leida} e índice \texttt{\{ destinatario: 1, leida: 1, createdAt: -1 \}} está especificado pero no implementado (Fase 7).
\end{itemize}

\chapter{Diagrama ER}
\label{ch:diagrama-er}

El campo \texttt{instrucciones} de la receta permanece como texto para conservar compatibilidad con los registros y la API existentes. La interfaz captura cada instrucción como un paso independiente y serializa los pasos en líneas numeradas (\texttt{Paso N. texto}); el detalle presenta las líneas como una lista ordenada. Las recetas históricas en un solo párrafo se conservan y se muestran como un único paso, sin dividir frases automáticamente.

\begin{figure}[H]
  \centering
  \includegraphics[width=0.72\textwidth]{build/diagrama-er.pdf}
  \caption{Diagrama entidad-relación del sistema.}
  \label{fig:diagrama-er}
\end{figure}

\chapter{Requerimientos no funcionales}
\label{ch:rnf}

El desarrollo sigue una arquitectura de tres capas: Frontend, Backend y Base de Datos. En la capa de base de datos se utiliza MongoDB como DBMS, con GridFS para archivos binarios. En la capa de Backend se usan servicios REST implementados en Node.js con Express y Mongoose, bajo el patrón Modelo Vista Controlador, consumidos desde la capa de Frontend. En la capa de Frontend se desarrollan componentes con Angular y TypeScript bajo el patrón de aplicación de página única (SPA), que se comunica con la base de datos a través de los servicios REST del Backend.

\section{Autenticación y autorización}

\begin{itemize}[leftmargin=*]
  \item El registro y el inicio de sesión se hacen con \path{POST /api/auth/register} y \path{POST /api/auth/login}.
  \item Las contraseñas se almacenan con hash bcrypt de 10 rondas; nunca en texto claro.
  \item La sesión se representa con un JWT de 30 días que viaja en la cabecera \texttt{Authorization: Bearer}.
  \item Todas las rutas salvo \texttt{/api/health}, las de autenticación y la lectura pública de imagen exigen un token válido; la ausencia o invalidez del token produce \texttt{401}.
  \item La autorización es por propiedad: solo el propietario puede editar, eliminar, cambiar la visibilidad o gestionar la imagen de su receta. Una operación sobre un recurso ajeno o inexistente produce \texttt{404}.
  \item El límite original de reintentos es de 10 intentos cada 15 minutos en registro e inicio de sesión.
  \item \textbf{Deuda conocida:} el cierre de sesión no invalida el token en el servidor; el \texttt{AuthGuard} del cliente comprueba la presencia del token, pero no su expiración antes de llamar al backend.
\end{itemize}

\section{Seguridad de datos}

\begin{itemize}[leftmargin=*]
  \item El servidor conserva la propiedad de \texttt{\_id}, \texttt{usuario}, contadores, reportes y métricas: nunca se aceptan desde el cliente.
  \item Se aplica saneado de entradas (\texttt{sanitizeFields}) en comentarios, reseñas, reportes y recetas.
  \item Se aplican límites de solicitudes: \texttt{apiLimiter} global de 100 solicitudes cada 15 minutos, \texttt{authLimiter} de 10 cada 15 minutos, \texttt{socialLimiter} de 30 por minuto y \texttt{createLimiter} de 10 por minuto.
  \item Los secretos (\texttt{MONGODB\_URI}, \texttt{JWT\_SECRET}) deben residir en variables de entorno y no en el repositorio.
  \item Se distinguen explícitamente los errores \texttt{400}, \texttt{401}, \texttt{404}, \texttt{429} y \texttt{500}.
  \item \textbf{Deuda conocida:} no se valida el \texttt{ObjectId} en todas las rutas (Mongoose lanza \texttt{CastError}); el saneado se aplica ruta a ruta en vez de globalmente; falta \texttt{app.set("trust proxy", 1)} para que el límite vea la IP real detrás de Caddy; \texttt{uploadLimiter} está definido pero no aplicado a la subida de imágenes; y hay secretos en el historial de \texttt{docker-compose.yml} pendientes de rotación.
\end{itemize}

\section{Almacenamiento}

\begin{itemize}[leftmargin=*]
  \item Los documentos residen en MongoDB 8 en Atlas y las imágenes en el bucket GridFS \texttt{uploads}.
  \item Solo se conserva una imagen vigente por receta; al reemplazar, quitar o borrar la receta se elimina el archivo anterior en GridFS.
  \item La optimización de imagen se hace en el navegador con \texttt{canvas}: lado mayor de 1600 px y JPEG de calidad 0.8 (medido: 2.4 MB a 186 KB). El servidor no usa \texttt{sharp}.
  \item El límite de carga es de 25 MB en el navegador y 5 MB en el servidor.
  \item La lectura de imagen es pública (\path{GET /api/recipes/image/:fileId}) y usa \texttt{Cache-Control: immutable}.
  \item \textbf{Deuda conocida:} no se verifican los bytes reales del archivo (magic bytes), solo el tipo MIME declarado; borrar una receta no elimina en cascada comentarios, reseñas ni actividad asociada.
\end{itemize}

\section{Integración}

\begin{itemize}[leftmargin=*]
  \item El cliente Angular consume la API REST bajo el prefijo \texttt{/api}, tanto por el proxy local (\path{proxy.conf.json}) como por Caddy en producción.
  \item Caddy central actúa como proxy inverso y terminación TLS para \path{https://foodloopc.duckdns.org}.
  \item Docker Compose construye y opera los contenedores \texttt{recetario-backend} y \texttt{recetario-frontend} sobre la red externa \texttt{caddy-central\_default}.
  \item El despliegue local usa procesos Node.js y Angular independientes (\texttt{npm run dev} y \texttt{npm start}).
  \item Los envíos de archivos usan \texttt{multipart/form-data} y \texttt{multer-gridfs-storage}.
\end{itemize}

\chapter{Planificación de ciclos de desarrollo}
\label{ch:planificacion}

\begin{center}
\begin{longtable}{p{1.9cm}p{8.4cm}p{3cm}}
  \toprule
  Ciclo & Qué abarca & Fechas \\
  \midrule
  \endfirsthead
  \toprule
  Ciclo & Qué abarca & Fechas \\
  \midrule
  \endhead
  Uno & Configuración del entorno, autenticación y CRUD de recetas (P01, P02). & [Completar: fecha] \\
  Dos & Imágenes en GridFS y visibilidad pública o privada (P03). & [Completar: fecha] \\
  Tres & Red social: interacciones, comentarios, reseñas, seguimiento y feed (P04, P05, P06). & [Completar: fecha] \\
  Cuatro & Notificaciones persistentes, perfil público y búsqueda global. & [Completar: fecha] \\
  \bottomrule
\end{longtable}
\end{center}

\chapter{Firmas de aceptación}
\label{ch:firmas}

Las personas que se indican a continuación declaran haber revisado el presente documento de especificaciones, requerimientos y criterios de aceptación de software del sistema \dercassistema\ y avalan su contenido.

\vspace{1cm}

\begin{center}
\begin{tabular}{p{7cm}p{5cm}}
  \toprule
  Nombre y Puesto & Firma y Sello \\
  \midrule
  \rule{0pt}{2.2cm}[Completar: contraparte informática] & \\[1.6cm]
  \rule{0pt}{2.2cm}[Completar: experto del dominio] & \\[1.6cm]
  \rule{0pt}{2.2cm}[Completar: responsable del despliegue] & \\[1.6cm]
  \rule{0pt}{2.2cm}[Completar: autor del documento] & \\[1.6cm]
  \bottomrule
\end{tabular}
\end{center}

\vspace{1cm}

\begin{center}
\begin{tabular}{p{7cm}}
  \toprule
  Lugar y fecha de aceptación \\
  \midrule
  \rule{0pt}{1.4cm} \\
  \bottomrule
\end{tabular}
\end{center}


\end{document}
